\documentclass[a4paper,12pt]{article}

\usepackage{fontsetup}
\usepackage[english, russian, ukrainian]{babel}

\usepackage{amsmath}

\newcommand\vect[1]{\vect{#1}}

\begin{document}

\subsection*{Детальний вивід середньої інтенсивності випромінювання для колового руху}

Розглянемо заряд $q$, що рухається по колу радіуса $r$ у площині $(x,y)$ зі сталою кутовою швидкістю $\omega = v/r$. Тоді:
\[
\vect{v}_q(t_q) = v\bigl( \sin(\omega t_q),\; \cos(\omega t_q),\; 0 \bigr),\qquad
\dot{\vect{v}}_q(t_q) = \omega v\bigl( \cos(\omega t_q),\; -\sin(\omega t_q),\; 0 \bigr),
\]
\[
\dot{v}^2 = \omega^2 v^2.
\]

Напрямок спостереження:
\[
\vect{n} = (\sin\theta,\; 0,\; \cos\theta),
\]
де $\theta$ --- кут між $\vect{n}$ і віссю $z$ (перпендикулярною до площини орбіти).

\subsubsection{Крок 1. Вираз для $Z_q$}
\[
\vect{v}_q\cdot\vect{n} = v\sin(\omega t_q)\sin\theta + v\cos(\omega t_q)\cdot 0 + 0\cdot\cos\theta
= v\sin\theta\sin(\omega t_q).
\]
Отже,
\[
Z_q = 1 - \frac{\vect{v}_q\cdot\vect{n}}{c} = 1 - \beta\sin\theta\sin\phi,\qquad
\beta = \frac{v}{c},\quad \phi = \omega t_q.
\]

\subsubsection{Крок 2. Миттєва інтенсивність $dI$ для $\vect{v}_q\perp\dot{\vect{v}}_q$}
З формули (4.1.14):
\[
dI = \frac{q^2}{4\pi c^3}\left[ \frac{\dot{v}^2}{Z_q^4} - \frac{(1-\beta^2)(\vect{n}\cdot\dot{\vect{v}}_q)^2}{Z_q^6} \right] do.
\]
Обчислимо $\vect{n}\cdot\dot{\vect{v}}_q$:
\[
\vect{n}\cdot\dot{\vect{v}}_q = \sin\theta\cdot(\omega v\cos\phi) + 0\cdot(-\omega v\sin\phi) + \cos\theta\cdot 0 = \omega v\sin\theta\cos\phi.
\]
Тому:
\[
(\vect{n}\cdot\dot{\vect{v}}_q)^2 = \omega^2 v^2\sin^2\theta\cos^2\phi.
\]
Підставляємо:
\[
dI = \frac{q^2}{4\pi c^3}\,\omega^2 v^2\left[ \frac{1}{Z_q^4} - \frac{(1-\beta^2)\sin^2\theta\cos^2\phi}{Z_q^6} \right] do.
\]

\subsubsection{Крок 3. Середнє за часом спостереження}
Середнє за великим періодом $T$:
\[
\langle dI \rangle = \frac{1}{T}\int_0^T dI(t)\,dt.
\]
Перехід до інтегрування по $t_q$:
\[
dt = \left(1 - \frac{\vect{v}_q\cdot\vect{n}}{c}\right) dt_q = Z_q\, dt_q.
\]
Оскільки $dI(t)$ залежить від $t$ через $t_q$, маємо:
\[
\langle dI \rangle = \frac{\omega}{2\pi} \int_0^{2\pi/\omega} dI(t_q)\, Z_q\, dt_q.
\]
Підставляємо $dI(t_q)$:
\[
\langle dI \rangle = \frac{\omega}{2\pi} \int_0^{2\pi/\omega} \frac{q^2}{4\pi c^3}\omega^2 v^2\left[ \frac{1}{Z_q^4} - \frac{(1-\beta^2)\sin^2\theta\cos^2\phi}{Z_q^6} \right] Z_q\, dt_q.
\]
Спрощуємо:
\[
\langle dI \rangle = \frac{q^2\omega^3 v^2}{8\pi^2 c^3} \int_0^{2\pi/\omega} \left[ \frac{1}{Z_q^3} - \frac{(1-\beta^2)\sin^2\theta\cos^2\phi}{Z_q^5} \right] dt_q.
\]
Заміна $dt_q = d\phi/\omega$, $\phi\in[0,2\pi]$:
\[
\langle dI \rangle = \frac{q^2\omega^2 v^2}{8\pi^2 c^3} \int_0^{2\pi} \left[ \frac{1}{Z_q^3} - \frac{(1-\beta^2)\sin^2\theta\cos^2\phi}{Z_q^5} \right] d\phi,
\]
де $Z_q = 1 - a\sin\phi$, $a = \beta\sin\theta$.

\subsubsection{Крок 4. Табличні інтеграли}
Відомі формули ($|a|<1$):
\[
J_3(a) = \int_0^{2\pi} \frac{d\phi}{(1-a\sin\phi)^3} = \frac{2\pi(2+a^2)}{(1-a^2)^{5/2}},
\]
\[
J_5(a) = \int_0^{2\pi} \frac{d\phi}{(1-a\sin\phi)^5} = \frac{2\pi(8+12a^2+3a^4)}{(1-a^2)^{9/2}}.
\]
Для інтеграла з $\cos^2\phi$:
\[
K_5(a) = \int_0^{2\pi} \frac{\cos^2\phi\,d\phi}{(1-a\sin\phi)^5} = \frac{\pi(4+3a^2)}{(1-a^2)^{7/2}}.
\]
(Цей результат можна отримати, виразивши $\cos^2\phi = 1-\sin^2\phi$ і використавши зв'язок між $J_n$ та похідними по $a$.)

Отже,
\[
\langle dI \rangle = \frac{q^2\omega^2 v^2}{8\pi^2 c^3}\Bigl[ J_3(a) - (1-\beta^2)\sin^2\theta\cdot K_5(a) \Bigr].
\]

\subsubsection{Крок 5. Підстановка $J_3$ та $K_5$}
\[
J_3(a) = \frac{2\pi(2+\beta^2\sin^2\theta)}{(1-\beta^2\sin^2\theta)^{5/2}},
\]
\[
K_5(a) = \frac{\pi(4+3\beta^2\sin^2\theta)}{(1-\beta^2\sin^2\theta)^{7/2}}.
\]
Тоді:
\[
\langle dI \rangle = \frac{q^2\omega^2 v^2}{8\pi^2 c^3}\left[
\frac{2\pi(2+\beta^2 s^2)}{(1-\beta^2 s^2)^{5/2}}
- (1-\beta^2)s^2\cdot \frac{\pi(4+3\beta^2 s^2)}{(1-\beta^2 s^2)^{7/2}}
\right],
\]
де $s\equiv\sin\theta$.

Виносимо $\pi$ за дужки:
\[
\langle dI \rangle = \frac{q^2\omega^2 v^2}{8\pi c^3}\left[
\frac{2(2+\beta^2 s^2)}{(1-\beta^2 s^2)^{5/2}}
- \frac{(1-\beta^2)s^2(4+3\beta^2 s^2)}{(1-\beta^2 s^2)^{7/2}}
\right].
\]

\subsubsection{Крок 6. Зведення до спільного знаменника}
Спільний знаменник $(1-\beta^2 s^2)^{7/2}$:
\[
\langle dI \rangle = \frac{q^2\omega^2 v^2}{8\pi c^3}\cdot
\frac{ 2(2+\beta^2 s^2)(1-\beta^2 s^2) - (1-\beta^2)s^2(4+3\beta^2 s^2) }
{(1-\beta^2 s^2)^{7/2}}.
\]

Обчислимо чисельник $N$:
\[
N = 2(2+\beta^2 s^2)(1-\beta^2 s^2) - (1-\beta^2)s^2(4+3\beta^2 s^2).
\]

Перший доданок:
\begin{multline*}
2(2+\beta^2 s^2)(1-\beta^2 s^2) = 2\bigl[2(1-\beta^2 s^2) + \beta^2 s^2(1-\beta^2 s^2)\bigr] = \\
= 2\bigl[2 - 2\beta^2 s^2 + \beta^2 s^2 - \beta^4 s^4\bigr] = \\
= 2\bigl[2 - \beta^2 s^2 - \beta^4 s^4\bigr] = 4 - 2\beta^2 s^2 - 2\beta^4 s^4.
\end{multline*}

Другий доданок:
\[
(1-\beta^2)s^2(4+3\beta^2 s^2) = s^2(4+3\beta^2 s^2) - \beta^2 s^2(4+3\beta^2 s^2)
= 4s^2 + 3\beta^2 s^4 - 4\beta^2 s^2 - 3\beta^4 s^4.
\]

Отже,
\[
N = \bigl(4 - 2\beta^2 s^2 - 2\beta^4 s^4\bigr) - \bigl(4s^2 + 3\beta^2 s^4 - 4\beta^2 s^2 - 3\beta^4 s^4\bigr).
\]

Збираємо подібні доданки:
\begin{align*}
\text{сталі: } & 4,\\
s^2: & -2\beta^2 s^2 - 4s^2 + 4\beta^2 s^2 = -4s^2 + 2\beta^2 s^2,\\
s^4: & -2\beta^4 s^4 - 3\beta^2 s^4 + 3\beta^4 s^4 = (\beta^4 - 3\beta^2) s^4.
\end{align*}
Таким чином,
\[
N = 4 - 4s^2 + 2\beta^2 s^2 - 3\beta^2 s^4 + \beta^4 s^4.
\]

\subsubsection{Крок 7. Перехід до $\cos\theta$ і спрощення}
Оскільки $s^2 = \sin^2\theta = 1 - \cos^2\theta$, підставимо:
\[
N = 4 - 4(1-\cos^2\theta) + 2\beta^2(1-\cos^2\theta) - 3\beta^2(1-\cos^2\theta)^2 + \beta^4(1-\cos^2\theta)^2.
\]

Спрощуємо перші два доданки: $4 - 4 + 4\cos^2\theta = 4\cos^2\theta$.

Далі:
\[
2\beta^2(1-\cos^2\theta) = 2\beta^2 - 2\beta^2\cos^2\theta.
\]

Розкриємо квадрати:
\[
(1-\cos^2\theta)^2 = 1 - 2\cos^2\theta + \cos^4\theta.
\]
Тоді:
\[
-3\beta^2(1-\cos^2\theta)^2 = -3\beta^2 + 6\beta^2\cos^2\theta - 3\beta^2\cos^4\theta,
\]
\[
+\beta^4(1-\cos^2\theta)^2 = \beta^4 - 2\beta^4\cos^2\theta + \beta^4\cos^4\theta.
\]

Додаємо все:
\begin{align*}
N = &\ 4\cos^2\theta \\
&+ 2\beta^2 - 2\beta^2\cos^2\theta \\
&- 3\beta^2 + 6\beta^2\cos^2\theta - 3\beta^2\cos^4\theta \\
&+ \beta^4 - 2\beta^4\cos^2\theta + \beta^4\cos^4\theta.
\end{align*}

Згрупуємо сталі: $2\beta^2 - 3\beta^2 + \beta^4 = -\beta^2 + \beta^4 = \beta^2(\beta^2-1) = -\beta^2(1-\beta^2)$.

Згрупуємо $\cos^2\theta$: $4\cos^2\theta - 2\beta^2\cos^2\theta + 6\beta^2\cos^2\theta - 2\beta^4\cos^2\theta
= 4\cos^2\theta + 4\beta^2\cos^2\theta - 2\beta^4\cos^2\theta
= 4\cos^2\theta(1+\beta^2) - 2\beta^4\cos^2\theta$.

Згрупуємо $\cos^4\theta$: $-3\beta^2\cos^4\theta + \beta^4\cos^4\theta = \beta^2(\beta^2-3)\cos^4\theta$.

Отже,
\[
N = -\beta^2(1-\beta^2) + \bigl[4(1+\beta^2)-2\beta^4\bigr]\cos^2\theta + \beta^2(\beta^2-3)\cos^4\theta.
\]

\subsubsection{Крок 8. Використання $\omega^2$ для магнітного поля}
Для колового руху в магнітному полі $B$:
\[
\omega = \frac{qB}{mc}\sqrt{1-\beta^2}, \qquad
\omega^2 v^2 = \frac{q^2 B^2}{m^2 c^2}(1-\beta^2)\cdot v^2 = \frac{q^2 B^2 v^2}{m^2 c^2}(1-\beta^2).
\]

Підставляємо у $\langle dI \rangle$:
\[
\langle dI \rangle = \frac{q^2}{8\pi c^3}\cdot \frac{q^2 B^2 v^2}{m^2 c^2}(1-\beta^2) \cdot \frac{N}{(1-\beta^2 s^2)^{7/2}}\, do.
\]

Тобто:
\[
\langle dI \rangle = \frac{q^4 B^2 v^2}{8\pi m^2 c^5}(1-\beta^2)\, \frac{N}{(1-\beta^2\sin^2\theta)^{7/2}}\, do.
\]

\subsubsection{Крок 9. Остаточний вигляд}
Після підстановки спрощеного $N$ отримуємо:
\[
\boxed{
\langle dI \rangle = \frac{q^4 B^2 v^2}{8\pi m^2 c^5}\left(1-\frac{v^2}{c^2}\right)
\frac{2 - \cos^2\theta - \frac{v^2}{4c^2}\left(1+3\frac{v^2}{c^2}\right)\cos^4\theta}
{\left(1 - \frac{v^2}{c^2}\cos^2\theta\right)^{7/2}}\; do.
}
\]

Це і є шуканий вираз після ``тривалих обчислень''.

\end{document}